\Needspace{10\baselineskip}
\item \textbf{Ajuste lineal y estimacion de un umbral de temperatura}.
\nopagebreak[4]

\begin{tcolorbox}[colback=blue!2!white,colframe=blue!55!black,title={Enunciado},
                  before skip=3pt,after skip=5pt,boxsep=3pt,top=3pt,bottom=3pt,
                  breakable]
\small
\raggedright
Durante el enfriamiento de una muestra se midio la temperatura \(T\), en grados
Celsius, en distintos tiempos \(t\), en minutos:
\[
t = [0,\ 2,\ 4,\ 6,\ 8,\ 10],
\]
\[
T = [82,\ 70,\ 59,\ 48,\ 38,\ 29].
\]

Suponga que, en este intervalo, la temperatura puede aproximarse por un modelo
lineal:
\[
T(t) = mt + b.
\]

Realice las siguientes operaciones:

\begin{itemize}
\setlength{\topsep}{4pt}
\setlength{\itemsep}{3pt}
\setlength{\parsep}{0pt}
\item[(a)] Importe \texttt{numpy} y \texttt{matplotlib.pyplot}. Defina los arreglos \texttt{t} y \texttt{T} usando \texttt{np.array}.
\item[(b)] Use \texttt{np.polyfit(t, T, deg=1)} para encontrar la pendiente \(m\) y el intercepto \(b\).
\item[(c)] Calcule los valores ajustados y guardelos en una variable llamada \texttt{T\_fit}.
\item[(d)] Calcule los residuos y el coeficiente \(R^2\). Para esto use:
\[
\texttt{residuos = T - T\_fit},
\]
\[
SS_{\mathrm{res}} = \sum_i \left(T_i - T_{\mathrm{fit},i}\right)^2,
\]
\[
SS_{\mathrm{tot}} = \sum_i \left(T_i - \overline{T}\right)^2,
\]
\[
R^2 = 1 - \frac{SS_{\mathrm{res}}}{SS_{\mathrm{tot}}}.
\]
\item[(e)] Grafique los datos medidos con \texttt{plt.scatter} y el ajuste lineal con \texttt{plt.plot}. Agregue etiquetas, leyenda y grilla.
\item[(f)] Use \texttt{np.roots} para estimar en que tiempo el modelo ajustado alcanza \(T=35^\circ\mathrm{C}\). Para ello resuelva:
\[
mt + b - 35 = 0.
\]
\item[(g)] Imprima \(m\), \(b\), \(R^2\) y el tiempo estimado para alcanzar \(35^\circ\mathrm{C}\). interprete brevemente dentro de un \texttt{print} el signo de \(m\).
\end{itemize}
\end{tcolorbox}

\begin{solution}
\begin{pythonjupyter}
# (a)
import numpy as np
import matplotlib.pyplot as plt

t = np.array([0, 2, 4, 6, 8, 10])
T = np.array([82, 70, 59, 48, 38, 29])

# (b)
m, b = np.polyfit(t, T, deg=1)

# (c)
T_fit = m*t + b

# (d)
residuos = T - T_fit
SS_res = np.sum(residuos**2)
SS_tot = np.sum((T - np.mean(T))**2)
R2 = 1 - SS_res/SS_tot

# (e)
plt.figure(figsize=(7, 5))
plt.scatter(t, T, label="datos medidos")
plt.plot(t, T_fit, "r-", label="ajuste lineal")
plt.xlabel("Tiempo (min)")
plt.ylabel("Temperatura (C)")
plt.legend()
plt.grid(True, alpha=0.3)
plt.show()

# (f)
coef_umbral = [m, b - 35]
raices = np.roots(coef_umbral)
tiempo_35 = raices[0].real

# (g)
print(f"m = {m:.3f} C/min")
print(f"b = {b:.3f} C")
print(f"R2 = {R2:.4f}")
print(f"Tiempo para T = 35 C: {tiempo_35:.3f} min")
print("La pendiente m es negativa, por lo tanto la temperatura disminuye")
print("a medida que avanza el tiempo.")
\end{pythonjupyter}
\end{solution}
